\section{Assembly, the explicit sequence, and limitations}
\label{sec:assembly}

Let \(\mu_{g,Q_g}\) be the empirical Frobenius measure of the pencil and let
\(m_g^{\mathrm{Haar}}\) be Haar probability measure on \(\USp(2g)\).
The preceding sections establish three comparisons:
\begin{align*}
 &\Pr(t_g\text{ is non-LI})=o(1),\\
 &d_{\mathrm{BL}}
 \bigl((T_{g,\lambda})_*\mu_{g,Q_g},
       (T_{g,\lambda})_*m_g^{\mathrm{Haar}}\bigr)=o(1),\\
 &T_{g,\lambda}(\mathsf U_g)\Longrightarrow\nu_\lambda(\PiSp)
 \quad(\mathsf U_g\text{ Haar}).
\end{align*}
Here the bounded--Lipschitz distance is taken on the Polish space
\((\Ptwo,\Wtwo)\).  Replacing the model law by the actual endpoint law changes
the expectation of every bounded--Lipschitz test by at most twice the first
probability.  The triangle inequality
therefore proves \eqref{eq:main-convergence}.

The limiting conditional law is almost surely absolutely continuous by
\Cref{prop:marked-stability}.  Hence evaluation on \((0,\infty)\) is almost
surely a continuity point for weak convergence of measures.  The continuous
mapping theorem gives \eqref{eq:density-convergence}.  Since these masses lie in \([0,1]\), their
expectations converge as well.  The two nondegeneracy clauses follow from
\Cref{prop:random-law-nondegenerate,prop:scalar-nondegenerate}.

This proves \cref{thm:main,cor:prime-race-densities}.  The explicit sequence in
\cref{ex:explicit-fields} is admissible by the calculation already recorded in
\eqref{eq:explicit-field-growth}.

\paragraph{Scope.}
The proof gives a conservative sufficient growth condition, not a practical or
optimal threshold.  More importantly, it does not establish the fixed-field
microscopic limit, nor does it prove an analogous theorem for the full
\(2g+1\)-parameter hyperelliptic ensemble.  The theorem is a quantitative
large-field/large-genus bridge for the explicit pencil \eqref{eq:pencil}.
